% erzeugt von 20_code/10_gearcore/scripts/thesis_appendix_tables.py, nicht von Hand ändern
\begin{table}[htbp]
    \caption{Geometrie der Prüfverzahnung kst-E, mit der Geometriekette aus der STplus-Eingabedatei des Falls gerechnet: Einzelräder, Werkzeuge, erzeugte Verzahnung und Prüfmaße. Die Zahndickenabmaße sind oberes und unteres Grenzabmaß, die Zahndicke und die Lückenweite gelten für die Nennverzahnung mit $x$, die Erzeugung für $x_{\mathrm{E}}$. Ein leeres Feld bedeutet, dass das Rad die Größe nicht hat.}
    \label{tab:anhang_geometrie_kste}
    \centering
    \footnotesize
    \setlength{\tabcolsep}{4pt}
    \begin{tabular}{L{6.2cm}lll}
        Größe & Einheit & Ritzel & Rad \\
        \hline
        \arrayrulecolor{black!30}
        $z$ &  & 51 & 52 \\
        \hline
        $m_{\mathrm{n}}$ & mm & $1{,}0000$ & $1{,}0000$ \\
        \hline
        $\alpha_{\mathrm{n}}$ & ° & $20{,}000$ & $20{,}000$ \\
        \hline
        $\beta$ & ° & $0{,}000$ & $0{,}000$ \\
        \hline
        $b$ & mm & $17{,}0000$ & $15{,}0000$ \\
        \hline
        $x$ &  & $0{,}2034$ & $0{,}3143$ \\
        \hline
        $E_{\mathrm{sns}} / E_{\mathrm{sni}}$ & µm & $-295{,}84$ / $-295{,}84$ & $-220{,}31$ / $-220{,}31$ \\
        \hline
        $x_{\mathrm{E}}$ &  & $-0{,}2030$ & $0{,}0117$ \\
        \hline
        $h_{\mathrm{aP0}}$ & mm & $1{,}1000$ & $1{,}2500$ \\
        \hline
        $\rho_{\mathrm{aP0}}$ & mm & $0{,}2000$ & $0{,}2000$ \\
        \hline
        $h_{\mathrm{FfP0}}$ & mm & $1{,}3000$ & $0{,}8456$ \\
        \hline
        $\alpha_{\mathrm{kP}}$ & ° &  & $45{,}000$ \\
        \hline
        $d$ & mm & $51{,}0000$ & $52{,}0000$ \\
        \hline
        $d_{\mathrm{b}}$ & mm & $47{,}9243$ & $48{,}8640$ \\
        \hline
        $d_{\mathrm{w}}$ & mm & $51{,}4951$ & $52{,}5049$ \\
        \hline
        $d_{\mathrm{a}}$ & mm & $52{,}8940$ & $54{,}0220$ \\
        \hline
        $d_{\mathrm{Fa}}$ & mm & $52{,}8940$ & $53{,}7876$ \\
        \hline
        $d_{\mathrm{Nf}}$ & mm & $50{,}3905$ & $51{,}3148$ \\
        \hline
        $d_{\mathrm{Ff}}$ & mm & $49{,}0811$ & $50{,}1566$ \\
        \hline
        $d_{\mathrm{fE}}$ & mm & $48{,}3940$ & $49{,}5234$ \\
        \hline
        $h$ & mm & $2{,}2500$ & $2{,}2493$ \\
        \hline
        $h_{\mathrm{a}}$ & mm & $0{,}9470$ & $1{,}0110$ \\
        \hline
        $h_{\mathrm{f}}$ & mm & $1{,}3030$ & $1{,}2383$ \\
        \hline
        $h_{\mathrm{K}}$ & mm & $0{,}0000$ & $0{,}1172$ \\
        \hline
        $s_{\mathrm{t}}$ & mm & $1{,}7189$ & $1{,}7996$ \\
        \hline
        $e_{\mathrm{t}}$ & mm & $1{,}4227$ & $1{,}3420$ \\
        \hline
        $s_{\mathrm{at}}$ & mm & $0{,}6715$ & $0{,}7766$ \\
        \hline
        $s_{\mathrm{aK}}$ & mm &  & $0{,}6342$ \\
        \hline
        $c$ & mm & $0{,}7913$ & $0{,}7920$ \\
        \hline
        $c_{\mathrm{F}}$ & mm & $0{,}6547$ & $0{,}5791$ \\
        \hline
        \multicolumn{4}{l}{Prüfmaße nach DIN 21773 mit $k$ und $D_{\mathrm{M}}$ nach der Regel von STplus} \\
        \hline
        $k$ &  & 6 & 6 \\
        \hline
        $W_{\mathrm{k}}$ (Nennmaß) & mm & $17{,}0901$ & $17{,}1800$ \\
        \hline
        $W_{\mathrm{k}}$ (oberes Grenzmaß) & mm & $16{,}8121$ & $16{,}9730$ \\
        \hline
        $W_{\mathrm{k}}$ (unteres Grenzmaß) & mm & $16{,}8121$ & $16{,}9730$ \\
        \hline
        $d_{\mathrm{M}}$ & mm & $50{,}7877$ & $51{,}7279$ \\
        \hline
        $D_{\mathrm{M}}$ & mm & $1{,}7500$ & $1{,}7500$ \\
        \hline
        $M_{\mathrm{dK}}$ (Nennmaß) & mm & $53{,}8464$ & $55{,}0640$ \\
        \hline
        $M_{\mathrm{dK}}$ (oberes Grenzmaß) & mm & $53{,}1099$ & $54{,}5321$ \\
        \hline
        $M_{\mathrm{dK}}$ (unteres Grenzmaß) & mm & $53{,}1099$ & $54{,}5321$ \\
        \hline
        $d_{\mathrm{K}}$ & mm & $51{,}3843$ & $52{,}7821$ \\
        \hline
        $s_{\mathrm{c}}$ (Nennmaß) & mm & $1{,}7185$ & $1{,}7993$ \\
        \hline
        $h_{\mathrm{c}}$ & mm & $0{,}9615$ & $1{,}0266$ \\
        \hline
        \arrayrulecolor{black}
    \end{tabular}
\end{table}
